% ============================================================================
% Fichier  : partie0/P00_le_matin_du_6_mars.tex
% Titre    : Prologue : le Matin du 6 Mars
% Auteur   : MINKA MI NGUIDJOÏ Thierry Emmanuel
% Version  : 1.0 -- novembre 2026
% ============================================================================
% Opération GAROUA vue par la direction. Mêmes faits techniques que le cours
% SEC3266 (connexion du 2 mars à 23 h 12, chiffrement du 5 mars à 21 h 40,
% découverte le 6 mars à 7 h 30). Cas fictif.
% ============================================================================

\chapter{Prologue : le Matin du 6 Mars}
\label{chap:prologue}
\chaptermark{Prologue}

\epigraph{%
    « Le plan n'est rien~; la planification est tout. »%
}{Formule attribuée à Dwight D.~Eisenhower}

\niveauChapitre{Fondamental}{Aucun~; le cours SEC3266 (Sécurité
informatique) éclaire les causes techniques.}

\begin{boxobjectifs}
\begin{itemize}
    \item Lire le déroulement d'une crise réelle du point de vue de la
          direction.
    \item Séparer ce qui a été constaté de ce qui a été supposé, et des
          opinions exprimées.
    \item Repérer, dans une crise, ce qui relevait du risque, de l'audit, de
          la continuité et de la décision.
    \item Formuler les questions auxquelles le cours répond.
\end{itemize}
\end{boxobjectifs}

\bigskip

% ============================================================================
\section*{Question directrice : qui décide, quand rien n'a été prévu~?}
% ============================================================================

Le vendredi 6~mars 2026, à 7~h~30, le responsable informatique de LiZbiZ
découvre que les serveurs de Douala et de Garoua sont chiffrés. Trois
heures plus tard, une cellule de crise improvisée se réunit dans le bureau
du directeur général. Personne n'a de liste des clients à prévenir, ni de
numéro de l'assureur, ni d'idée du temps qu'il faudra pour facturer à
nouveau. Ce prologue raconte cette journée et les suivantes, sans les
juger~; les chapitres qui suivent donnent les outils pour les juger.

\medskip
\textbf{Plan du chapitre.}
\begin{enumerate}
    \item LiZbiZ avant la crise (section~\ref{sec:p00-avant}).
    \item La journée du 6 mars (section~\ref{sec:p00-journee}).
    \item Les jours suivants (section~\ref{sec:p00-suite}).
    \item Constatations, suppositions, opinions (section~\ref{sec:p00-lire}).
    \item Les questions du cours (section~\ref{sec:p00-questions}).
\end{enumerate}

\separateur

% ============================================================================
\section{LiZbiZ avant la Crise}
\label{sec:p00-avant}
% ============================================================================

LiZbiZ\footnote{LiZbiZ, ses dirigeants, ses clients et ses prestataires
sont fictifs, comme les dates, les montants et les adresses du cours~; les
méthodes et les normes citées sont réelles.} est une société anonyme de
logistique dont le siège est à Douala. Elle emploie 150~salariés et
réalise un chiffre d'affaires d'environ 24~milliards de FCFA par an, soit
quelque 80~millions de FCFA par jour ouvré. Elle exploite quatre centres
de données, à Douala, Yaoundé, Garoua et Bafoussam, et a confié en
septembre 2025 l'exploitation de Garoua et de Douala au prestataire SAHEL
INFOGÉRANCE.

\begin{tableaucours}{L{4.0cm}L{8.2cm}}{%
    \tableauHeader{Acteur} &
    \tableauHeader{Rôle pendant la crise}
}
M.~ETEME & Directeur général~; préside la cellule de crise \\
Mme~MBALLA & Directrice administrative et financière~; trésorerie,
assurance, clients \\
M.~KOUNOU & Responsable informatique~; découvre l'attaque, pilote la
reprise \\
Mme~NKENG & Contrôleuse interne~; tient la main courante à partir du
7~mars \\
M.~NGUELE & Directeur des opérations~; expéditions, agences, mode dégradé \\
Mme~ABOMO & Chargée de communication~; clients, presse, personnel \\
M.~NANGA & Chef d'agence de Garoua \\
Me~TSALA & Avocat de la société~; plainte, rançon, données personnelles \\
\end{tableaucours}
\vspace{-0.8em}
\begin{center}{\small\itshape Tableau~0.1. Les acteurs de l'Opération GAROUA}\end{center}

\textbf{Constatation.} En février 2026, LiZbiZ n'a ni registre des
risques, ni plan de continuité, ni plan de gestion de crise. Le dernier
audit de sécurité date de 2023 et ne couvrait pas le prestataire. Le
contrat avec SAHEL INFOGÉRANCE ne mentionne ni délai de notification d'un
incident, ni obligation de second facteur.

% ============================================================================
\section{La Journée du 6 Mars}
\label{sec:p00-journee}
% ============================================================================

La chronologie ci-dessous a été reconstituée par Mme~NKENG à partir des
messages, des appels et des souvenirs des participants~; aucune main
courante n'a été tenue avant le 7~mars.

\begin{tableaucours}{L{2.2cm}L{10.0cm}}{%
    \tableauHeader{Heure} &
    \tableauHeader{Événement}
}
07~h~30 & M.~KOUNOU constate que les serveurs de Douala et de Garoua sont
chiffrés~; un message de rançon réclame 85\,000~dollars en cryptoactifs. \\
07~h~45 & Il appelle M.~ETEME, qui est en route pour Yaoundé. \\
08~h~40 & M.~NGUELE apprend par un client que le portail de suivi ne
répond plus. \\
09~h~05 & Les liaisons entre les sites sont toujours ouvertes~; personne
n'a autorité pour les couper sans le directeur général. \\
10~h~30 & Réunion dans le bureau du DG, par téléphone pour M.~ETEME~:
huit personnes, aucun ordre du jour, aucun compte rendu. \\
11~h~10 & Décision de couper les liaisons entre les sites. Yaoundé et
Bafoussam sont épargnés. \\
13~h~00 & Mme~MBALLA cherche le contrat d'assurance~: la police ne couvre
pas les cyberattaques. \\
15~h~30 & Les agents des agences reçoivent la consigne orale de noter les
expéditions sur papier. \\
18~h~30 & Une photographie du message de rançon, prise par un salarié,
circule sur les réseaux sociaux. \\
20~h~00 & Un journaliste appelle le standard~; personne n'est désigné pour
répondre. \\
\end{tableaucours}
\vspace{-0.8em}
\begin{center}{\small\itshape Tableau~0.2. La journée du 6 mars, reconstituée}\end{center}

\begin{boxscenario}{Opération GAROUA. Trois heures et quarante minutes}
\textbf{Constatation.} Entre la découverte (7~h~30) et la coupure des
liaisons (11~h~10), 3~h~40 se sont écoulées. Pendant ce temps, aucune
action de confinement n'a été prise, faute de personne autorisée à la
décider.

\medskip
\textbf{Analyse} (confiance moyenne). Les serveurs de Yaoundé et de
Bafoussam n'ont pas été touchés~; il n'est donc pas établi que le délai ait
aggravé les dommages. Il est établi, en revanche, que l'organisation ne
savait pas qui pouvait décider~: c'est le premier défaut de gestion de
crise (Chapitre~\ref{chap:cellule}).
\end{boxscenario}

% ============================================================================
\section{Les Jours Suivants}
\label{sec:p00-suite}
% ============================================================================

\begin{tableaucours}{L{2.2cm}L{10.0cm}}{%
    \tableauHeader{Date} &
    \tableauHeader{Événement}
}
7~mars & Première réunion avec main courante~; communiqué aux clients à
9~h~; plainte déposée auprès de la police judiciaire à Douala. \\
8~mars & Après avis de Me~TSALA, M.~ETEME décide de ne pas payer la
rançon. \\
9~mars & Les mots de passe d'administration sont changés. \\
12~mars & Premier comité de sécurité~: registre des risques établi après
coup. \\
16~mars & Reprise de la facturation, onze jours après l'attaque. \\
19~mars & L'attaquant publie sur son site un échantillon de 2~Go de
données clients~; seconde vague d'appels. \\
2~avril & Réunion de retour d'expérience, avec la direction et SAHEL
INFOGÉRANCE. \\
\end{tableaucours}
\vspace{-0.8em}
\begin{center}{\small\itshape Tableau~0.3. Les jours suivants}\end{center}

Le 2~avril, Mme~MBALLA présente une première estimation des coûts de la
crise.

\begin{tableaucours}{L{7.4cm}R{2.8cm}}{%
    \tableauHeader{Poste} &
    \tableauHeader{Millions de FCFA}
}
Remise en état (prestataires, matériel, licences) & 96 \\
Ressaisie de six semaines de factures & 18 \\
Pénalités contractuelles versées aux clients & 74 \\
Clients perdus (marge annuelle estimée) & 110 \\
Conseil juridique et communication & 14 \\
\textbf{Total} & \textbf{312} \\
\end{tableaucours}
\vspace{-0.8em}
\begin{center}{\small\itshape Tableau~0.4. Les coûts de l'Opération GAROUA (estimation du 2 avril)}\end{center}

\begin{boxalert}
\textbf{Violation typique.} Présenter « 312~millions » comme un fait.

\smallskip
\textbf{Règle~:} trois postes sont constatés (remise en état, ressaisie,
pénalités, soit 188~millions)~; deux sont estimés (clients perdus,
conseil). Une estimation se présente avec ses hypothèses et sa marge
(Chapitre~\ref{chap:quantitatif}).
\end{boxalert}

% ============================================================================
\section{Constatations, Suppositions, Opinions}
\label{sec:p00-lire}
% ============================================================================

Au retour d'expérience du 2~avril, chacun a sa lecture de la crise. Le
cours demande de séparer trois registres.

\begin{tableaucours}{L{3.0cm}L{9.2cm}}{%
    \tableauHeader{Registre} &
    \tableauHeader{Exemples entendus le 2 avril}
}
Constatation & « Les liaisons ont été coupées à 11~h~10 »~; « la police
d'assurance exclut les cyberattaques » \\
Supposition & « Si l'on avait coupé à 8~h, rien n'aurait changé »~;
« les clients perdus le sont à cause de la photo » \\
Opinion & « Le prestataire est incompétent »~; « on a bien géré, vu les
circonstances » \\
\end{tableaucours}
\vspace{-0.8em}
\begin{center}{\small\itshape Tableau~0.5. Trois registres}\end{center}

\begin{boxmetier}{Directeur de crise}
\textbf{Faire tenir la main courante dès la première minute.} Le directeur
de crise désigne, avant toute décision, une personne qui consigne l'heure,
le fait, la source et la décision. Sans elle, la crise se raconte après
coup, de mémoire.
\end{boxmetier}

\begin{boxmetier}{Chargé de communication de crise}
\textbf{Parler avant les réseaux sociaux.} Un message bref et exact, publié
tôt, vaut mieux qu'un communiqué complet publié après la rumeur. À 18~h~30,
la rumeur a parlé la première.
\end{boxmetier}

% ============================================================================
\section{Les Questions du Cours}
\label{sec:p00-questions}
% ============================================================================

\begin{tableaucours}{L{3.4cm}L{6.6cm}L{2.2cm}}{%
    \tableauHeader{Partie} &
    \tableauHeader{Question posée par la crise} &
    \tableauHeader{Chapitres}
}
I. Fondements & Pourquoi personne n'avait-il jugé ce risque prioritaire~?
& \ref{chap:risque} à~\ref{chap:droit} \\
II. Analyser & Quelle perte fallait-il craindre, et avec quelle
probabilité~? & \ref{chap:identifier} à~\ref{chap:traiter} \\
III. Auditer & Un audit aurait-il vu le compte sans second facteur~? &
\ref{chap:audit} à~\ref{chap:rapport} \\
IV. Continuer & Pourquoi onze jours pour facturer~? & \ref{chap:incidents}
à~\ref{chap:exercer} \\
V. Crise & Qui devait décider à 7~h~45, et que fallait-il dire~? &
\ref{chap:crise} à~\ref{chap:exercices-crise} \\
VI. Apprendre & Qu'a-t-on appris, et comment le mesurer~? &
\ref{chap:retex} \\
\end{tableaucours}
\vspace{-0.8em}
\begin{center}{\small\itshape Tableau~0.6. Les questions du cours}\end{center}

\separateur

\begin{boxresume}
\textbf{Ce qu'il faut retenir de ce chapitre~:}
\begin{itemize}
    \item \textbf{La crise de LiZbiZ} n'a pas commencé le 6~mars~: elle
          tenait à l'absence de registre, de plan et d'audit.
    \item \textbf{Le premier défaut} du 6~mars fut de ne pas savoir qui
          décidait~: 3~h~40 sans confinement.
    \item \textbf{Sans main courante}, une crise se reconstitue de mémoire.
    \item \textbf{Les coûts} mêlent constatations et estimations~; il faut
          les distinguer.
    \item \textbf{Constatation, supposition, opinion}~: trois registres à
          ne jamais confondre.
\end{itemize}
\end{boxresume}

\begin{boxinfo}
\textbf{Points de contrôle de l'élève.} Je sais~: retracer la chronologie
du 6~mars~; nommer les acteurs et leur rôle~; distinguer les coûts
constatés des coûts estimés~; classer une affirmation en constatation,
supposition ou opinion~; relier chaque défaut à une partie du cours.
\end{boxinfo}

\bigskip

\begin{boxexo}[title={\ding{45}\quad Exercice~0.1. Le délai \niveaubase}]
Calculez le délai entre la découverte de l'attaque et la décision de
couper les liaisons, puis entre la découverte et le premier communiqué aux
clients.
\end{boxexo}

\begin{boxexo}[title={\ding{45}\quad Exercice~0.2. Les registres \niveaubase}]
Classez en constatation, supposition ou opinion~: « la rançon aurait été
payée sans Me~TSALA »~; « la facturation a repris le 16~mars »~; « SAHEL
INFOGÉRANCE aurait dû prévenir »~; « le portail ne répondait plus à
8~h~40 ».
\end{boxexo}

\begin{boxexo}[title={\ding{45}\quad Exercice~0.3. Les coûts \niveaubase}]
Dans le Tableau~0.4, quelle part du total est constatée, et quelle part
est estimée~? Donnez les montants et les pourcentages.
\end{boxexo}

\begin{boxexo}[title={\ding{45}\quad Exercice~0.4. Les absents \niveaumoyen}]
Quelles fonctions manquaient à la réunion de 10~h~30 pour décider vite~?
Quelles informations auraient dû être disponibles à l'ouverture de la
réunion~?
\end{boxexo}

\begin{boxexo}[title={\ding{45}\quad Exercice~0.5. Le message \niveaumoyen}]
Rédigez le message de trois phrases que LiZbiZ aurait pu publier à
19~h le 6~mars, après la diffusion de la photographie.
\end{boxexo}

\begin{boxexo}[title={\ding{45}\quad Exercice~0.6. L'assurance \niveaumoyen}]
La police d'assurance excluait les cyberattaques. Quelle étape de la
gestion des risques aurait dû le révéler, et à quel moment~?
\end{boxexo}

\begin{boxexo}[title={\ding{45}\quad Exercice~0.7. La rançon \niveaumoyen}]
Listez les éléments que la cellule de crise devait réunir avant de
décider de payer ou non la rançon.
\end{boxexo}

\begin{boxexo}[title={\ding{45}\quad Exercice~0.8. La seconde vague \niveaumoyen}]
Le 19~mars, un échantillon de données est publié. Quelles parties
prenantes faut-il informer, et dans quel ordre~?
\end{boxexo}

\begin{boxexo}[title={\ding{45}\quad Exercice~0.9. La chronologie idéale \niveauexpert}]
Réécrivez la journée du 6~mars, de 7~h~30 à 20~h, telle qu'elle aurait dû
se dérouler avec une cellule de crise préparée.
\end{boxexo}

\begin{boxexo}[title={\ding{45}\quad Exercice~0.10. Le programme \niveauexpert}]
À partir du Tableau~0.6, proposez à M.~ETEME un programme de six mois
pour que la prochaine crise soit mieux gérée~: actions, responsables,
échéances.
\end{boxexo}

\bigskip

\begin{boxinfo}
\textbf{Pour aller plus loin.}
\begin{itemize}
    \item P.~Lagadec, \emph{La gestion des crises}~\cite{lagadec1991}.
    \item ISO~22361, lignes directrices pour la gestion de
          crise~\cite{iso22361_2022}.
\end{itemize}
\end{boxinfo}

\begin{boxinfo}
\textbf{Transition.} Avant de gérer une crise, il faut savoir ce qu'est un
risque. Le Chapitre~\ref{chap:risque} distingue risque et incertitude, et
montre ce qu'est une bonne décision sous incertitude.
\end{boxinfo}
